
%% ===================================================================
%% LP-134 canonical naming (2025-12-15):
%% Per LP-134 (Lux Chain Topology), T-Chain has been split into:
%%   M-Chain — MPC ceremonies for bridge custody of EXTERNAL wallets
%%             (CGGMP21, FROST, Pulsar-general)
%%   F-Chain — FHE compute + TFHE bootstrap-key generation (encrypted EVM)
%% The legacy name "T-Chain" is retained ONLY for `teleportvm` (LP-6332).
%% Where this paper says "T-Chain MPC" / "T-Chain threshold governance" /
%% "T-Chain TFHE" / "T-Chain custody", read as M-Chain (MPC) or F-Chain (FHE).
%% ===================================================================

\documentclass[11pt]{article}
\usepackage{amsmath,amssymb,amsthm}
\usepackage{booktabs}
\usepackage{enumitem}
\usepackage{hyperref}
\usepackage[margin=1in]{geometry}

\hypersetup{colorlinks=true,linkcolor=black,citecolor=blue,urlcolor=blue}

\newtheorem{theorem}{Theorem}
\newtheorem{lemma}[theorem]{Lemma}
\newtheorem{corollary}[theorem]{Corollary}
\newtheorem{definition}{Definition}
\newtheorem{remark}{Remark}

% Theorem numbers in this document are deliberately offset so that the
% main results are 7.5, 7.6, 7.7, 7.8, matching the LP-020 paper numbering.
% (Theorems 7.1--7.4 live in LP-020 itself.)
\setcounter{theorem}{4}

\title{QuasarCert Soundness:\\
{\large Canonical Formal Specification, Quasar 3.0}}
\author{Zach Kelling\\
\textit{Lux Industries, Inc.}\\
\texttt{research@lux.network}}
\date{Spec freeze: 2025-12-15 \\ Activation: 2025-12-25}

\begin{document}
\maketitle

\begin{abstract}
This note is the canonical formal specification of QuasarCert soundness
for the Quasar 3.0 chain activation on 2025-12-25.  QuasarCert is a
three-lane finality certificate
$\mathcal{C} = (\sigma_{\mathrm{BLS}}, \sigma_{\mathrm{RT}},
\pi_{\mathrm{ZK}})$ with each lane bound to a common
\emph{certificate subject} that fixes chain id, epoch, round, mode and
the per-chain root commitments $(P, Q, Z, A, B, M, F)$.  We prove four
results with theorem numbers aligned to LP-020:
(7.5) classical soundness, requiring an adversary to forge all three
lanes against three independent hardness assumptions;
(7.6) parallel-lane liveness, allowing the three lanes to be aggregated
independently and in parallel;
(7.7) post-quantum safety, the standalone PQ predicate
$\mathsf{VerifyPQ} = \mathsf{RT.Verify} \land \mathsf{ZK.Verify}$;
(7.8) subject-binding replay protection, ruling out cross-epoch reuse of
any valid cert.  We also address the rejection of threshold ML-DSA in
3.0 (App.~A), the GPU verifier in LP-132 (\S7), and the concrete
post-quantum parameter analysis (App.~E).

This proof establishes soundness of the QuasarCert wire format defined in LP-020, with constituent layers specified in LP-070 (ML-DSA), LP-073 (Pulsar; Ringtail-derived lattice threshold kernel with key-era lifecycle), and LP-075 (BLS aggregate).

\paragraph{Naming evolution (2026-03).} LP-073 was originally titled ``Ringtail.'' It now describes \textbf{Pulsar}, the Lux production fork at \texttt{github.com/luxfi/pulsar}, which inherits Ringtail's Sign1/Sign2/Combine math byte-equal but adds a key-era lifecycle (LSS-managed generations, verifiable resharing, activation certificate, rollback) suitable for permissionless validator-set rotation. Throughout this proof, ``Ringtail'' refers to the lattice signing primitive (the Sign1/Sign2/Combine math); ``Pulsar'' refers to the systems-layer composition that wraps it for dynamic permissionless deployment. Soundness results below derive from the SUF-CMA reduction of the underlying Ringtail signing scheme; lifecycle additions (resharing, activation, rollback) are orthogonal to the per-signature unforgeability argument and are addressed separately in LP-073's lifecycle section.
\end{abstract}

%% ==================================================================
\section{QuasarCert}
\label{sec:cert}
%% ==================================================================

\begin{definition}[Validator set]
\label{def:val}
$\mathcal{V} = \{v_1, \ldots, v_n\}$, $|\mathcal{V}| = n$.  Stake
distribution $\mathrm{stk}: \mathcal{V} \to \mathbb{N}$ with total
stake $S = \sum_i \mathrm{stk}(v_i)$.  At most a stake fraction $f < 1/3$
is Byzantine.  Threshold $\tau = 2/3$, i.e.\ a valid quorum collects
stake $\geq \lceil 2S/3 \rceil$.
\end{definition}

\begin{definition}[Certificate subject -- Def.~6.1 of LP-020]
\label{def:subject}
The \emph{certificate subject} is the byte string
\[
  \mathsf{subj} =
  \mathsf{chain\_id} \,\|\, \mathsf{epoch} \,\|\, \mathsf{round}
  \,\|\, \mathsf{mode}
  \,\|\, P \,\|\, Q \,\|\, Z \,\|\, A \,\|\, B \,\|\, M \,\|\, F
  \,\|\, \mathsf{parent\_*\_hashes}
  \,\|\, \mathsf{gas/fee}
\]
where $P$ = \texttt{pchain\_validator\_root}, $Q$ =
\texttt{qchain\_ceremony\_root}, $Z$ = \texttt{zchain\_vk\_root},
$A$ = \texttt{achain\_state\_root}, $B$ = \texttt{bchain\_state\_root},
$M$ = \texttt{mchain\_msg\_root}, $F$ = \texttt{fchain\_state\_root}.
All three lanes sign or prove against $\mathsf{subj}$.
\end{definition}

\begin{definition}[QuasarCert -- Def.~6.2 of LP-020]
\label{def:cert}
For subject $\mathsf{subj}$, a \emph{QuasarCert} is the triple
\[
  \mathcal{C} = (\sigma_{\mathrm{BLS}}, \sigma_{\mathrm{RT}},
                  \pi_{\mathrm{ZK}})
\]
where:
\begin{itemize}[nosep]
  \item $\sigma_{\mathrm{BLS}}$: BLS12-381 aggregate from validators of
    stake $\geq \tau S$ on $\mathsf{subj}$, verified against
    $\mathsf{pchain\_validator\_root}$;
  \item $\sigma_{\mathrm{RT}}$: Ringtail two-round threshold signature
    aggregating shares of stake $\geq \tau S$ on $\mathsf{subj}$,
    verified against $\mathsf{qchain\_ceremony\_root}$;
  \item $\pi_{\mathrm{ZK}}$: Z-Chain Groth16 proof that $n$ ML-DSA-65
    signatures on $\mathsf{subj}$ verify under the public-key vector
    committed by $\mathsf{zchain\_vk\_root}$.
\end{itemize}
$\mathsf{Verify}(\mathcal{C}) = 1$ iff every lane verifies and every
lane binds the \emph{same} $\mathsf{subj}$.
\end{definition}

\begin{definition}[Hardness assumptions]
\label{def:assumptions}
We isolate one assumption per lane.
\begin{enumerate}[nosep]
  \item \textbf{DL/co-CDH on BLS12-381} (classical).  Discrete log /
    co-CDH in $(\mathbb{G}_1, \mathbb{G}_2, \mathbb{G}_T)$.  BLS
    aggregate signatures with proof-of-possession are EUF-CMA under
    co-CDH.
  \item \textbf{Module-LWE} (post-quantum).  $\mathsf{MLWE}_{q,k,\ell,
    \chi}$ for the Ringtail / ML-DSA-65 parameter sets in App.~E.
  \item \textbf{Module-SIS} (post-quantum).  $\mathsf{MSIS}_{q,k,\beta}$
    for the same parameter sets.  Combined with MLWE, gives EUF-CMA for
    ML-DSA-65 (FIPS~204) and SUF-CMA for Ringtail.
  \item \textbf{KEA / AGM} (auxiliary).  Groth16 is knowledge-sound in
    the algebraic group model on BLS12-381.
\end{enumerate}
\end{definition}

%% ==================================================================
\section{Adversary Model}
\label{sec:adv}
%% ==================================================================

\begin{definition}[Adversary]
\label{def:adv}
A non-uniform polynomial-time adversary $\mathcal{A}$ controls a
subset $F \subset \mathcal{V}$ of validators with stake share $< \tau$.
$\mathcal{A}$ is given:
\begin{itemize}[nosep]
  \item the public keys of all validators (BLS, Ringtail, ML-DSA);
  \item signing oracles for the corrupt set $F$;
  \item the public verification key $\mathsf{vk}$ for the Z-Chain
    Groth16 circuit;
  \item the random oracle $H$ (programmable in the proof).
\end{itemize}
We consider two variants.
\begin{description}[nosep]
  \item[$\mathcal{A}_C$ -- classical PPT.]  Standard probabilistic
    polynomial time.
  \item[$\mathcal{A}_Q$ -- quantum PPT.]  Polynomial-time on a quantum
    Turing machine; can run Shor and Grover, has full quantum random
    oracle access (QROM).
\end{description}
$\mathcal{A}$ \emph{wins} the soundness game if it outputs
$(\mathsf{subj}^\star, \mathcal{C}^\star)$ such that
$\mathsf{Verify}(\mathcal{C}^\star) = 1$ and $\mathsf{subj}^\star$ was
not endorsed by any honest stake-weighted quorum.
\end{definition}

\begin{remark}[Static vs adaptive]
The proofs in Sections~\ref{sec:soundness}--\ref{sec:replay} are stated
under \emph{static} corruption.  Adaptive corruption is handled in
App.~C via secure-erasure or programmable-RO hybrids, with at most an
$O(n \cdot Q_H)$ multiplicative loss.  All Lux validators perform
Go-side memory zeroing of Ringtail signing randomness; HSM-backed
T-Chain signing on the live network provides hardware-enforced
erasures, so Cor.~\ref{cor:adaptive-erasure} applies tightly.
\end{remark}

%% ==================================================================
\section{Soundness}
\label{sec:soundness}
%% ==================================================================

\begin{theorem}[QuasarCert classical soundness]
\label{thm:main}
Under the assumptions of Def.~\ref{def:assumptions}, if
$\mathcal{C}^\star = (\sigma_{\mathrm{BLS}}^\star,
\sigma_{\mathrm{RT}}^\star, \pi_{\mathrm{ZK}}^\star)$ verifies against a
subject $\mathsf{subj}^\star$ that was not endorsed by validators of
stake $\geq \tau S$, then $\mathcal{A}_C$'s advantage satisfies
\[
  \mathsf{Adv}^{\mathrm{forge}}_{\mathcal{A}_C}
  \;\leq\;
  \min\!\bigl(
    \epsilon_{\mathrm{BLS}}^{\mathrm{coCDH}},\;
    \epsilon_{\mathrm{RT}}^{\mathrm{MLWE,MSIS}},\;
    \epsilon_{\mathrm{ZK}}^{\mathrm{KEA}} +
    n \cdot \epsilon_{\mathrm{MLDSA}}^{\mathrm{MLWE,MSIS}}
  \bigr)
  \;=\; \mathsf{negl}(\lambda).
\]
\end{theorem}

\begin{proof}
We argue per lane and then take the conjunction.

\medskip\noindent
\textbf{Lane 1 (BLS).}  $\sigma_{\mathrm{BLS}}^\star$ is a BLS12-381
aggregate over $\mathsf{subj}^\star$.  Proof-of-possession on each
validator's $\mathsf{pchain\_validator\_root}$ entry rules out
rogue-key attacks~\cite{boneh2001short}.  By EUF-CMA of BLS aggregate
signatures under co-CDH, a valid $\sigma_{\mathrm{BLS}}^\star$ implies
that each validator $v_i$ in the aggregator set
$S_{\mathrm{BLS}}^\star$ produced an individual BLS signature on
$\mathsf{subj}^\star$.  Stake-weighting then gives
$\mathrm{stk}(S_{\mathrm{BLS}}^\star) \geq \tau S$.  Forging without
$\tau$ stake costs $\epsilon_{\mathrm{BLS}}^{\mathrm{coCDH}}$.

\medskip\noindent
\textbf{Lane 2 (Ringtail).}  $\sigma_{\mathrm{RT}}^\star$ is the
output of the two-round Ringtail signing protocol bound to
$\mathsf{qchain\_ceremony\_root}$ (i.e., the verification key produced
by the Q-Chain DKG).  By the SUF-CMA proof of Ringtail under MLWE/MSIS,
a valid $\sigma_{\mathrm{RT}}^\star$ on $\mathsf{subj}^\star$ implies a
share set $S_{\mathrm{RT}}^\star$ with
$\mathrm{stk}(S_{\mathrm{RT}}^\star) \geq \tau S$.  Forging costs at
most $\epsilon_{\mathrm{RT}}^{\mathrm{MLWE,MSIS}}$.

\medskip\noindent
\textbf{Lane 3 (Z-Chain Groth16 of ML-DSA).}  $\pi_{\mathrm{ZK}}^\star$
proves that for the public-key vector committed by
$\mathsf{zchain\_vk\_root}$, there exist valid ML-DSA-65 signatures
$(\sigma_1^{\mathrm{ML}}, \ldots, \sigma_n^{\mathrm{ML}})$ on
$\mathsf{subj}^\star$.  Knowledge-soundness in the AGM yields an
extractor $\mathcal{E}$ that recovers these signatures with
overwhelming probability: any deviation has advantage
$\epsilon_{\mathrm{ZK}}^{\mathrm{KEA}}$.  Each extracted ML-DSA-65
signature is itself EUF-CMA under MLWE/MSIS, so for any uncorrupted
$v_i$ a valid extracted signature without $v_i$'s key costs
$\epsilon_{\mathrm{MLDSA}}^{\mathrm{MLWE,MSIS}}$ (with up to $n$ guessing
loss for which validator was framed).

\medskip\noindent
\textbf{Conjunction.}  $\mathsf{Verify}$ is the AND of the three
predicates.  An attacker who breaks fewer than three lanes still has
at least one lane verifying only with negligible probability.  Hence
$\mathsf{Adv}^{\mathrm{forge}}_{\mathcal{A}_C}$ is upper-bounded by the
\emph{minimum} of the per-lane advantages, which is negligible.

\medskip\noindent
\textbf{Cross-lane independence.}  The three bounds rely on three
disjoint problems: discrete log on BLS12-381, MLWE, MSIS.  No known
reduction transfers a break of one to a break of another at the
parameter sets of Def.~\ref{def:assumptions}; this is exactly the
hedging argument of pq-finality-no-bls.tex \S3.2.  Compromise of any
one of $\{\mathsf{pchain\_validator\_root}, \mathsf{qchain\_ceremony\_root},
\mathsf{zchain\_vk\_root}\}$ disables only the corresponding lane.
\end{proof}

\begin{corollary}[No conflicting certificates]
\label{cor:no-conflict}
If $\mathcal{C}, \mathcal{C}'$ both verify on subjects $\mathsf{subj},
\mathsf{subj}'$ that share $(\mathsf{chain\_id}, \mathsf{epoch},
\mathsf{round})$ but disagree on the chain-state roots, then
$> \tau$ stake is Byzantine.
\end{corollary}

\begin{proof}
By Lane~1, $\mathrm{stk}(S_{\mathrm{BLS}}) \geq \tau S$ and
$\mathrm{stk}(S_{\mathrm{BLS}}') \geq \tau S$.  Stake-weighted quorum
intersection gives
$\mathrm{stk}(S_{\mathrm{BLS}} \cap S_{\mathrm{BLS}}') \geq (2\tau - 1) S
= S/3$.  Any validator in the intersection signed two distinct subjects
at the same $(\mathsf{chain\_id}, \mathsf{epoch}, \mathsf{round})$,
which is detectable equivocation; if Byzantine stake is $< S/3$ this
is impossible.
\end{proof}

%% ==================================================================
\section{Parallel Lane Liveness}
\label{sec:parallel}
%% ==================================================================

\begin{theorem}[Parallel-lane liveness]
\label{thm:parallel}
The three lanes of QuasarCert can be aggregated independently and in
parallel.  Specifically, let $T_{\mathrm{BLS}}, T_{\mathrm{RT}},
T_{\mathrm{ZK}}$ be the wall-clock durations of the three lane
aggregations under partial synchrony with bounded delay $\Delta$ and
honest stake $\geq \tau$.  Then there exists an honest schedule that
produces $\mathcal{C}$ in time
\[
  T_{\mathrm{cert}}
  \;\leq\;
  \max\!\bigl(T_{\mathrm{BLS}},\; T_{\mathrm{RT}},\; T_{\mathrm{ZK}}\bigr)
  + \Delta.
\]
\end{theorem}

\begin{proof}[Proof sketch]
The three aggregations communicate only through the shared
$\mathsf{subj}$.  Once $\mathsf{subj}$ is fixed for a round (which
requires only the chain-state roots, all available at round start),
each lane runs a separate aggregation network:
BLS aggregator on the P-Chain validator network;
Ringtail two-round protocol on the Q-Chain DKG mesh;
Z-Chain Groth16 prover (epoch-amortised, asynchronous to block
production).  Honest validators contribute to all lanes they hold keys
for; under stake $\geq \tau$ each lane independently completes within
its own bound.  The certificate is finalised when the last lane lands;
the extra $\Delta$ accounts for one round of network propagation to
combine the three artifacts.

\medskip\noindent
The Z-Chain proof, in particular, is generated by the Z-Chain prover
service over an epoch window and lands asynchronously: the epoch's
$\pi_{\mathrm{ZK}}$ becomes part of every QuasarCert in that epoch
once posted.  This is the LP-020 Section~6.4 ``epoch ZK landing''
mechanism and does not block per-block finality of the BLS+Ringtail
core.
\end{proof}

%% ==================================================================
\section{Post-Quantum Safety}
\label{sec:pq}
%% ==================================================================

\begin{theorem}[Post-quantum safety]
\label{thm:pq}
Define the standalone PQ predicate
$\mathsf{VerifyPQ}(\mathcal{C}) = \mathsf{RT.Verify}(\sigma_{\mathrm{RT}})
\land \mathsf{ZK.Verify}(\pi_{\mathrm{ZK}})$.  Against a quantum
adversary $\mathcal{A}_Q$,
\[
  \mathsf{Adv}^{\mathsf{VerifyPQ}}_{\mathcal{A}_Q}
  \;\leq\;
  2 \cdot \max\!\bigl(
    \epsilon_{\mathrm{MLWE}}^{Q},\;
    \epsilon_{\mathrm{MSIS}}^{Q}
  \bigr)
  \;=\; \mathsf{negl}(\lambda).
\]
That is: even if $\mathcal{A}_Q$ trivially breaks BLS via Shor, the
PQ-only predicate remains sound.
\end{theorem}

\begin{proof}
$\mathcal{A}_Q$ must produce $\sigma_{\mathrm{RT}}^\star$ that verifies
against $\mathsf{qchain\_ceremony\_root}$ and $\pi_{\mathrm{ZK}}^\star$
that verifies against $\mathsf{zchain\_vk\_root}$.

\medskip\noindent
\textit{Forging $\sigma_{\mathrm{RT}}$ quantumly} requires breaking
MLWE or MSIS at the Ringtail dimension $d = 2048$, modulus
$\log_2 q = 48$ (App.~E).  The best known quantum cost is BDGL sieving
with Grover speedup, giving $T_{\mathrm{quantum}} \geq 2^{130}$
operations.

\medskip\noindent
\textit{Forging $\pi_{\mathrm{ZK}}$ quantumly} is delicate: Groth16's
classical pairing-based knowledge-soundness reduces to discrete log on
BLS12-381, which is broken by Shor.  In the post-quantum setting we
fall back to the soundness of the underlying ML-DSA-65 signatures it
proves.  A quantum forger of $\pi_{\mathrm{ZK}}$ either (a) extracts
no real witness and yet the proof verifies (breaks knowledge-soundness
in the AGM, classical only -- not assumed against $\mathcal{A}_Q$), or
(b) extracts ML-DSA-65 signatures the adversary did not legitimately
hold.  Case (b) costs $\epsilon_{\mathrm{MLDSA}}^{Q} \geq 2^{128}$ at
ML-DSA-65 (NIST Cat.~3) per signature, with at most $n$ guessing loss.
The standalone-PQ guarantee for $\pi_{\mathrm{ZK}}$ thus reduces to the
PQ unforgeability of ML-DSA-65 signatures committed by
$\mathsf{zchain\_vk\_root}$.

\medskip\noindent
A guessing reduction picks the targeted lane with probability $1/2$ and
inherits the corresponding bound; combined with $\mathsf{negl}$ per
lane this gives the stated $\mathsf{Adv}$.

\medskip\noindent
\textbf{Future-work caveat.}  Replacing Groth16 with a STARK-of-pairing
or hash-based SNARK eliminates the dependence on classical
knowledge-soundness in $\pi_{\mathrm{ZK}}$.  This is tracked as the
``post-quantum SNARK swap'' in LP-020 \S11 and is out of scope for 3.0
final.
\end{proof}

\begin{remark}[Combined three-lane PQ guarantee]
The full QuasarCert remains sound against $\mathcal{A}_Q$ as long as
\emph{at least one} of $\{$Ringtail, ML-DSA path$\}$ survives.
Equivalently: the survival predicate is ``1 of 2 PQ pillars''.  This
matches the pq-finality-no-bls.tex Proposition (Hedged PQ security)
under the conjunction ``MLWE hard $\lor$ MSIS hard,'' and is the
operative claim for the 2025-12-25 activation.
\end{remark}

%% ==================================================================
\section{Subject-Binding Replay Protection}
\label{sec:replay}
%% ==================================================================

This section is new in 3.0 final.  The result formalises why the
\emph{certificate subject} (Def.~\ref{def:subject}) makes a valid
QuasarCert from epoch $e_1$ unusable at any other epoch
$e_2 \neq e_1$.

\begin{theorem}[Subject-binding replay protection]
\label{thm:replay}
Let $\mathcal{C}$ be a QuasarCert that verifies for subject
$\mathsf{subj}_1$ at epoch $e_1$.  Let $\mathsf{subj}_2$ differ from
$\mathsf{subj}_1$ in any field of Def.~\ref{def:subject}.  Then the
probability that $\mathcal{C}$ also verifies for $\mathsf{subj}_2$
under the assumptions of Def.~\ref{def:assumptions} is
$\mathsf{negl}(\lambda)$.  In particular, the cert cannot be replayed
across distinct $\mathsf{epoch}$ values (or distinct
$\mathsf{chain\_id}$, $\mathsf{round}$, $\mathsf{mode}$, or any chain
root in the subject).
\end{theorem}

\begin{proof}
We show that each lane is collision-resistant on the subject input.

\medskip\noindent
\textbf{BLS.}  $\sigma_{\mathrm{BLS}}$ is computed as
$\prod_i H_{\mathbb{G}_2}(\mathsf{subj})^{sk_i}$ where
$H_{\mathbb{G}_2}$ is a hash-to-curve.  If
$\sigma_{\mathrm{BLS}}$ verifies for both $\mathsf{subj}_1$ and
$\mathsf{subj}_2 \neq \mathsf{subj}_1$, then there exists
$\sigma_{\mathrm{BLS}}$ such that
$e(\sigma_{\mathrm{BLS}}, g_1) = e(H_{\mathbb{G}_2}(\mathsf{subj}_1),
\mathsf{apk}) = e(H_{\mathbb{G}_2}(\mathsf{subj}_2), \mathsf{apk})$.
The pairing equality forces
$H_{\mathbb{G}_2}(\mathsf{subj}_1) = H_{\mathbb{G}_2}(\mathsf{subj}_2)$
(in $\mathbb{G}_2$), which is a hash-to-curve collision and hence
negligible.

\medskip\noindent
\textbf{Ringtail.}  The Ringtail challenge is
$c = H(\mathbf{w}, \mathsf{subj})$.  If $\sigma_{\mathrm{RT}}$ verifies
for both subjects with the same commitment $\mathbf{w}$, then
$H(\mathbf{w}, \mathsf{subj}_1) = H(\mathbf{w}, \mathsf{subj}_2)$,
again a hash collision.  If the commitments differ, this is a fresh
forgery and reduces to MLWE/MSIS.

\medskip\noindent
\textbf{Groth16.}  $\pi_{\mathrm{ZK}}$ encodes $\mathsf{subj}$ as a
public input bound at SRS preprocessing.  Verifying the same proof
against a different public input with non-negligible probability
contradicts knowledge-soundness in the AGM (the verifier equation
$e(A, B) = e(\alpha, \beta) \cdot e(\sum_i \mathsf{subj}[i] \cdot
\mathsf{IC}_i, \gamma) \cdot e(C, \delta)$ is linear in
$\mathsf{subj}$ and pinning a single proof to two distinct linear
combinations is a collision in the AGM extractor).

\medskip\noindent
Each lane's replay probability is therefore $\mathsf{negl}(\lambda)$;
since $\mathsf{Verify}$ is the AND of the three, replay across
distinct subjects has at most the minimum of these advantages, which
remains negligible.
\end{proof}

\begin{remark}[Concrete fields blocked]
Because $\mathsf{subj}$ contains $\mathsf{epoch}$, $\mathsf{round}$,
and $\mathsf{chain\_id}$ explicitly, replay across epochs, across
rounds within an epoch, or across chains (e.g., devnet to mainnet) is
ruled out individually.  The chain-root fields
$(P, Q, Z, A, B, M, F)$ further block replay against any state
divergence: even if two distinct epochs share epoch number (impossible
by P-Chain rules but treated formally), the $A/B/F$ roots disagree as
soon as any inclusion-list transaction differs.
\end{remark}

%% ==================================================================
\section{GPU Verifier Refinement (Informal)}
\label{sec:gpu}
%% ==================================================================

LP-132 specifies the QuasarGPU adapter.  Its core entry point is
\texttt{drain\_cert\_lane(lane, batch)} which produces a verdict
$\{\mathsf{accept}, \mathsf{reject}, \mathsf{abstain}\}$ for each cert
in $\mathsf{batch}$.  The CPU reference is a sequential loop calling
the per-lane verifier.

\paragraph{Cross-backend equivalence (LP-132 invariant).}
Let $V_{\mathrm{cpu}}$ and $V_{\mathrm{gpu}}$ denote the verdict
functions.  LP-132 requires
\[
  V_{\mathrm{gpu}}(\pi(\mathsf{batch})) = \pi(V_{\mathrm{cpu}}(\mathsf{batch}))
  \quad\text{for any canonical-tx-index permutation } \pi.
\]
That is, the GPU implementation differs from CPU only by deterministic
reordering.  The reordering is induced by per-warp scheduling, which
the adapter resolves at output by sorting on the canonical tx index.

\paragraph{Argument.}
Each lane verifier is a pure function on
$(\mathsf{subj}, \text{lane artifact})$ with no shared state across
batch elements.  GPU parallelism replicates this function across
threads with thread-local memory; the only global state read is the
verification key roots (read-only).  Hence each per-element verdict is
identical to the CPU verdict.  The output ordering is canonicalised by
sorting on the tx index, eliminating the only nondeterminism.  By
inspection of the LP-132 reference loop, the GPU adapter produces
verdict-equivalent output up to permutation, and the canonical sort
removes the permutation.

\paragraph{What this is not.}
This is an \emph{informal} argument, not a machine-checked proof.  A
machine-checked refinement would require reasoning about CUDA / HIP
memory models and is tracked in the \texttt{lean/Consensus/} GPU
refinement task.  The cross-backend differential test in LP-132
\texttt{tests/cross\_backend\_test.go} provides empirical evidence at
the engineering boundary.

%% ==================================================================
\section{Pulsar Key-Era Lifecycle Lemmas}
\label{sec:lifecycle}
%% ==================================================================

This section formalises the lifecycle properties of LP-073 Pulsar that
sit \emph{above} the per-signature SUF-CMA argument. The Ringtail kernel
of Pulsar is treated as an oracle that produces a signature
$\sigma_{\mathrm{RT}}$ verifying under
$\mathsf{GK} = (\mathbf{A}, \tilde{\mathbf{b}})$ given $\geq t$ shares
$\{\mathbf{s}_i\}$ that satisfy the Lagrange-recombination identity for
the master secret $\mathbf{s}$. We prove four lemmas that together
formalise the ``persistent body, rotating beam'' property of a key era.

\begin{definition}[Pulsar share state]
\label{def:share-state}
A Pulsar share state at validator $v$ is a tuple
\[
  \pi_v = \bigl(
    \eta,\; g,\; r,\;
    V,\; t,\;
    \mathbf{A},\; \tilde{\mathbf{b}},\;
    \mathbf{s}_v,\; \boldsymbol\lambda_v,\;
    \mathsf{seeds}_v,\; \mathsf{macKeys}_v
  \bigr)
\]
where $\eta \in \mathbb{N}$ is the \emph{KeyEraID}, $g \in \mathbb{N}$
the \emph{Generation}, $r \in \mathbb{N}$ the \emph{RollbackFrom}
counter ($r = 0$ on forward transitions),
$V \subseteq \mathcal{V}$ the active validator subset of size $|V| = K$,
$t$ the threshold, $\mathbf{A} \in R_q^{m_{\mathrm{pk}} \times
n_{\mathrm{vec}}}$ and $\tilde{\mathbf{b}} \in R_\xi^{m_{\mathrm{pk}}}$
the persistent group-key components, $\mathbf{s}_v$ the share assigned
to $v$, $\boldsymbol\lambda_v$ the precomputed Lagrange coefficient for
$v$ in $V$, and $\mathsf{seeds}_v, \mathsf{macKeys}_v$ the pairwise PRF
seeds and BLAKE3 MAC keys for $v$ vs.\ every $w \in V \setminus \{v\}$.
\end{definition}

\begin{definition}[GroupKey hash]
\label{def:gkh}
The \emph{GroupKey hash} $\mathsf{GKH}(\pi_v)$ is the 32-byte digest
\[
  \mathsf{GKH}(\pi_v) = \mathsf{BLAKE3}\bigl(
    \mathsf{WriteTo}(\mathbf{A}) \,\|\,
    \mathsf{WriteTo}(\tilde{\mathbf{b}})
  \bigr)[:32].
\]
$\mathsf{WriteTo}$ is the LP-073 \S4 wire serialisation: little-endian
\texttt{uint64} coefficient stream with 8/16-byte length prefixes for
$\mathsf{Vector}$/$\mathsf{Matrix}$. Two share states $\pi_v, \pi_w$
\emph{share a GroupKey} iff $\mathsf{GKH}(\pi_v) = \mathsf{GKH}(\pi_w)$.
\end{definition}

\begin{definition}[Lifecycle transitions]
\label{def:transitions}
Let $\Pi$ denote a global Pulsar state $(\pi_v)_{v \in V}$. We
distinguish three transitions:
\begin{itemize}[nosep]
  \item \textbf{Refresh} ($\Pi \xrightarrow{\mathrm{Refresh}} \Pi'$):
    HJKY97 zero-polynomial proactive refresh, same $V, t$, same
    $(\mathbf{A}, \tilde{\mathbf{b}})$, share distribution rotates.
  \item \textbf{Reshare} ($\Pi \xrightarrow{\mathrm{Reshare}(V', t')} \Pi'$):
    Desmedt--Jajodia97 redistribution to validator subset $V'$ with
    threshold $t'$; same $(\mathbf{A}, \tilde{\mathbf{b}})$.
  \item \textbf{Reanchor} ($\Pi \xrightarrow{\mathrm{Reanchor}} \Pi'$):
    fresh ceremony or DKG produces a new $(\mathbf{A}', \tilde{\mathbf{b}}')$
    and bumps $\eta$.
\end{itemize}
\end{definition}

\begin{lemma}[Generation monotonicity]
\label{lem:gen-mono}
For every Pulsar transition $\Pi \to \Pi'$ at fixed $\eta$ and any
$v \in V \cap V'$, $g(\pi'_v) = g(\pi_v) + 1$.
\end{lemma}

\begin{proof}
The LSS coordinator (LP-073~\S9 (LSS-Pulsar)) advances $g$ by one
on each successful Refresh or Reshare. Failed activation triggers
\textsc{Rollback}; we treat \textsc{Rollback}$(g_{\mathrm{target}})$ as
a separate transition $\Pi \to \Pi'$ for which the post-state has
$g(\pi'_v) = g(\pi_v) + 1$ \emph{and} $r(\pi'_v) = g_{\mathrm{target}}$.
This preserves monotonicity --- $g$ is bumped on every transition,
including reverts --- while marking lineage via $r$. The activation
certificate (Theorem~\ref{thm:activation-soundness} below) carries
$g$ as an explicit field, so any verifier can audit monotonicity from
the on-chain record. The Lean encoding of this property appears as
\texttt{reshare\_threshold\_maintained} and the implicit chain of
generations in \texttt{Crypto.Threshold.Reshare} of the canonical
proof tree at \texttt{lean/Crypto/Threshold/Reshare.lean}.
\end{proof}

\begin{lemma}[KeyEraID stability across Reshare]
\label{lem:eta-stable}
For every $\Pi \xrightarrow{\mathrm{Reshare}(V', t')} \Pi'$ and every
$v \in V'$, $\eta(\pi'_v) = \eta(\pi_v)$ for any $v \in V \cap V'$
(or $\eta(\pi_v)$ inherited from the launch transition for new
$v \in V' \setminus V$). $\eta$ bumps only at Reanchor.
\end{lemma}

\begin{proof}
Reshare carries $(\mathbf{A}, \tilde{\mathbf{b}})$ across by
construction (Desmedt--Jajodia identity, Lemma~\ref{lem:gk-identical}
below); $\eta$ is the chain-state pointer to the GroupKey lineage and
is updated only when the underlying $(\mathbf{A}, \tilde{\mathbf{b}})$
changes, i.e., at Reanchor. The LSS adapter contract
(LP-073~\S9 (LSS-Pulsar).3) explicitly preserves $\eta$ across
the call \texttt{DynamicResharePulsar}; the LP-073 acceptance test
\texttt{TestKeyEraIDPreservedAcrossReshare}
(\texttt{threshold/protocols/lss/lss\_pulsar\_test.go}) verifies this
property at the bytes-on-disk level for the Go canonical.
\end{proof}

\begin{lemma}[GroupKey byte-identity across Refresh and Reshare]
\label{lem:gk-identical}
For every $\Pi \to \Pi'$ that is either Refresh or
$\mathrm{Reshare}(V', t')$, and any $v \in V'$,
\[
  \mathsf{GKH}(\pi'_v) = \mathsf{GKH}(\pi_w)
  \quad \forall w \in V \cap V'.
\]
That is, the BLAKE3 digest of the wire-serialised
$(\mathbf{A}, \tilde{\mathbf{b}})$ is identical pre- and post-transition,
not merely structurally equal.
\end{lemma}

\begin{proof}
Refresh and Reshare update only the share polynomial; neither alters
the public matrix $\mathbf{A}$ nor the rounded public-key vector
$\tilde{\mathbf{b}}$. We verify the Lagrange-recombination identity
that underwrites this preservation. Let $f(x) \in R_q[X]$ be the
old-set sharing polynomial with $f(0) = \mathbf{s}$ and
$f(\alpha_v) = \mathbf{s}_v$ for $v \in V$. After
$\mathrm{Reshare}(V', t')$, the new polynomial is constructed as
\[
  G(x) = \sum_{i \in Q} \lambda^Q_i(0) \cdot g_i(x),
  \qquad g_i(0) = \mathbf{s}_i,\; \deg g_i = t' - 1,
\]
where $Q \subseteq V$ is a qualified old subset of size
$\geq t_{\mathrm{old}}$ and $\lambda^Q_i(0)$ are the Lagrange
coefficients of $Q$ evaluated at $0$. Then
\begin{align}
  G(0) \;=\; \sum_{i \in Q} \lambda^Q_i(0) \cdot g_i(0)
       \;=\; \sum_{i \in Q} \lambda^Q_i(0) \cdot \mathbf{s}_i
       \;=\; \sum_{i \in Q} \lambda^Q_i(0) \cdot f(\alpha_i)
       \;=\; f(0)
       \;=\; \mathbf{s}, \label{eq:lagrange}
\end{align}
where the penultimate equality is the standard Lagrange interpolation
identity over $R_q$, applied componentwise to the polynomial
coefficients. Hence the master secret $\mathbf{s}$ is preserved, the
LWE error $\mathbf{e}$ is unchanged (it was sampled at genesis to form
$\tilde{\mathbf{b}}$ and erased with the dealer state; signers only
need shares of $\mathbf{s}$, never $\mathbf{e}$), and the public-key
vector $\tilde{\mathbf{b}} = \lfloor \mathbf{A} \mathbf{s} +
\mathbf{e} \rceil_\xi$ remains unchanged because both
$(\mathbf{A}, \mathbf{s}, \mathbf{e})$ remain unchanged.
$\mathsf{WriteTo}$ is a deterministic byte function of
$(\mathbf{A}, \tilde{\mathbf{b}})$, so the BLAKE3 digest is identical.
For Refresh, an analogous argument applies with $z_i(x)$ a
zero-polynomial of degree $t-1$ \cite{hjky97}: the constant term
contributes nothing to $G(0)$, so $\mathbf{s}$, $\mathbf{e}$, and
hence $\tilde{\mathbf{b}}$ are byte-preserved.

The identity~\eqref{eq:lagrange} is mechanised in
\texttt{lean/Crypto/Threshold/Reshare.lean} as the axiom
\texttt{reshare\_same\_secret} together with theorem
\texttt{reshare\_preserves\_public\_key}, which derives
$\mathsf{pubkey\_from\_shares}(\mathsf{newShares}, t')
= \mathsf{pubkey}(\mathsf{key})$ from \texttt{reshare\_same\_secret}.
\end{proof}

\begin{lemma}[Rollback-lineage forward reconstruction]
\label{lem:rollback-lineage}
Let
$\Pi_0 \xrightarrow{\mathrm{Reshare}} \Pi_1 \xrightarrow{\mathrm{Reshare}}
\cdots \xrightarrow{\mathrm{Reshare}} \Pi_k$ be a chain of forward
transitions at fixed $\eta$, and suppose
$\Pi_k \xrightarrow{\mathrm{Rollback}(j)} \Pi_{k+1}$ for some $j < k$.
Then the chain
$(\eta_i, g_i, r_i)_{i=0}^{k+1}$ uniquely reconstructs the
forward-only path $\Pi_0 \to \Pi_1 \to \cdots \to \Pi_j$ that the
rollback restored to.
\end{lemma}

\begin{proof}
Each transition records $(\eta_i, g_i, r_i)$ on chain. By
Lemma~\ref{lem:gen-mono}, $g_i$ is strictly monotone in $i$. By
construction, $r_i = 0$ iff transition $i$ is forward; $r_i = j > 0$
iff transition $i$ is a rollback to generation $j$. To reconstruct the
forward-only path: walk the chain in reverse, skipping any transition
$i$ with $r_i > 0$ together with all transitions $i' > j$ such that
$g_{i'} > j$ and $r_{i'} = 0$ for $i' \leq i$ (these are the reverted
forward steps). The remaining sequence is the canonical forward
lineage of any state restored to. Uniqueness follows from the
strict monotonicity of $g$ within an era (Lemma~\ref{lem:gen-mono})
and the fact that $r$ unambiguously labels rollback targets.
\end{proof}

\begin{theorem}[Activation-certificate soundness]
\label{thm:activation-soundness}
Let $\mathsf{ActivationMsg}_{\eta, g+1}$ be the wire payload of
LP-073~\S8.3 (activation cert), binding the new generation
$g+1$ at fixed era $\eta$. The verifier
$\mathsf{VerifyActivation}(\Pi, \Pi', \sigma)$ accepts iff
$\sigma$ is a valid Pulsar threshold signature on
$\mathsf{ActivationMsg}_{\eta, g+1}$ that verifies under the
\emph{unchanged} GroupKey $(\mathbf{A}, \tilde{\mathbf{b}})$ of $\Pi$.
Then under MLWE/MSIS at the LP-073 \S2 parameters,
\[
  \Pr\bigl[
    \mathsf{VerifyActivation} = 1 \wedge
    \pi'_v \text{ does not satisfy~\eqref{eq:lagrange}}
  \bigr] \leq \epsilon_{\mathrm{RT}}^{\mathrm{MLWE,MSIS}} = \mathsf{negl}(\lambda).
\]
That is: $\mathsf{VerifyActivation}$ accepts iff the new committee's
shares can sign under the unchanged GroupKey, which by SUF-CMA reduces
to the new shares satisfying the Lagrange-recombination
identity~\eqref{eq:lagrange}.
\end{theorem}

\begin{proof}
Forward direction. If $\Pi'$ satisfies~\eqref{eq:lagrange}, then by
Lemma~\ref{lem:gk-identical} the GroupKey is byte-identical, and the
honest threshold signing protocol (Ringtail Round 1 / Round 2) on the
new shares produces a $\sigma$ that verifies under the same
$(\mathbf{A}, \tilde{\mathbf{b}})$. $\mathsf{VerifyActivation}$
accepts.

Reverse direction. If $\mathsf{VerifyActivation}$ accepts, then by the
SUF-CMA reduction of Ringtail under MLWE/MSIS (Theorem~5.1 of
LP-073 below, i.e.\ Theorem~\ref{thm:main} restricted to the Ringtail
lane), there exists a share set $\{\mathbf{s}'_v\}_{v \in V'}$ of stake
$\geq t'$ that produced the signature, and these shares jointly
recombine to the master secret $\mathbf{s}$ via Lagrange interpolation.
By the verifier equation $A z - b c$ and the $L_2$ bound, the recovered
$\mathbf{s}$ is consistent with $(\mathbf{A}, \tilde{\mathbf{b}})$ ---
otherwise the forger has produced a fresh MLWE/MSIS instance, an event
of probability at most $\epsilon_{\mathrm{RT}}^{\mathrm{MLWE,MSIS}}$.
Hence the new shares satisfy~\eqref{eq:lagrange} except with negligible
probability.
\end{proof}

\begin{remark}[What activation does \emph{not} prove]
$\mathsf{VerifyActivation}$ is necessary but not sufficient for full
Byzantine attribution: it certifies that the new committee can sign,
not which old or new party (if any) sent a malformed contribution.
Attribution is the responsibility of the complaint and disqualification
pipeline (lattice Pedersen commit verification, LP-073~\S8.5
(Pedersen DKG over $R_q$)). Failed activation triggers the non-punitive
rollback ladder of LP-073~\S8.4 (rollback); safety and
liveness do not depend on slashing-based attribution.
\end{remark}

%% ==================================================================
\section{Gaps and Future Work}
\label{sec:gaps}
%% ==================================================================

The 3.0 final scope leaves three gaps explicitly open for future work,
each tracked in LP-020.

\begin{enumerate}[nosep]
  \item \textbf{Threshold ML-DSA.}  Not used in 3.0 due to a
    rejection-sampling circularity (see App.~A).  Resolved-or-rejected
    once the Celi--del~Pino--Espitau--Niot--Prest construction is
    stable.
  \item \textbf{Adaptive corruption without erasures.}  Tightness of
    the $O(n \cdot Q_H)$ loss in App.~C; potential gain from a
    Fischlin-style transform.
  \item \textbf{Post-quantum SNARK.}  Replacing Groth16 with a STARK-
    or hash-based SNARK to eliminate the conditional on classical
    knowledge-soundness in Theorem~\ref{thm:pq}.
\end{enumerate}


%% ==================================================================
%% BIBLIOGRAPHY
%% ==================================================================
\begin{thebibliography}{20}

\bibitem{boneh2001short} D.~Boneh, B.~Lynn, H.~Shacham, ``Short
  Signatures from the Weil Pairing,'' \emph{ASIACRYPT}, 2001.

\bibitem{ringtail} Z.~Kelling, V.~Seesahai, ``Ringtail: Lattice-Based
  PQ Threshold Signatures for Blockchain Consensus,''
  Lux Industries, 2024.

\bibitem{nist-pqc} NIST, ``FIPS 204: ML-DSA (Module-Lattice-Based
  Digital Signature Algorithm),'' 2024.

\bibitem{groth2016} J.~Groth, ``On the Size of Pairing-Based
  Non-Interactive Arguments,'' \emph{EUROCRYPT}, 2016.

\bibitem{regev2005} O.~Regev, ``On Lattices, Learning with Errors,
  Random Linear Codes, and Cryptography,'' \emph{J.~ACM}, 56(6), 2009.

\bibitem{canetti2001} R.~Canetti, ``Universally Composable Security:
  A New Paradigm for Cryptographic Protocols,'' \emph{FOCS}, 2001.

\bibitem{garay2015} J.~Garay, A.~Kiayias, N.~Leonardos, ``The Bitcoin
  Backbone Protocol: Analysis and Applications,'' \emph{EUROCRYPT},
  2015.

\bibitem{fischlin2005} M.~Fischlin, ``Communication-Efficient
  Non-Interactive Proofs of Knowledge with Online Extractors,''
  \emph{CRYPTO}, 2005.

\bibitem{bowe2017} S.~Bowe, A.~Gabizon, I.~Miers, ``Scalable
  Multi-Party Computation for zk-SNARK Parameters in the Random Beacon
  Model,'' IACR ePrint 2017/1050, 2017.

\bibitem{plonk2019} A.~Gabizon, Z.~J.~Williamson, O.~Ciobotaru,
  ``PLONK: Permutations over Lagrange-Bases for Oecumenical
  Noninteractive Arguments of Knowledge,'' IACR ePrint 2019/953.

\bibitem{marlin2020} A.~Chiesa, Y.~Hu, M.~Maller, P.~Mishra,
  N.~Vesely, N.~Ward, ``Marlin: Preprocessing zkSNARKs with Universal
  and Updatable SRS,'' \emph{EUROCRYPT}, 2020.

\bibitem{bdgl2016} A.~Becker, L.~Ducas, N.~Gama, T.~Laarhoven,
  ``New Directions in Nearest Neighbor Searching with Applications to
  Lattice Sieving,'' \emph{SODA}, 2016.

\bibitem{nist-pqc-cost} NIST, ``Status Report on the Third Round of
  the NIST Post-Quantum Cryptography Standardization Process,''
  NISTIR~8413, 2022.

\bibitem{espitau2024} T.~Espitau, M.~Tibouchi, A.~Wallet, ``Practical
  Fiat-Shamir with Aborts Threshold Signatures from Lattices,''
  IACR ePrint, 2024.

\bibitem{laarhoven2015} T.~Laarhoven, ``Sieving for Shortest Vectors
  in Lattices Using Angular Locality-Sensitive Hashing,''
  \emph{CRYPTO}, 2015.

\bibitem{albrecht-estimator} M.~R.~Albrecht, R.~Player, S.~Scott,
  ``On the Concrete Hardness of Learning with Errors,''
  \emph{J.~Math.~Cryptol.}, 9(3):169--203, 2015.

\bibitem{nist-mldsa-spec} NIST, ``FIPS 204: Module-Lattice-Based
  Digital Signature Algorithm (ML-DSA),'' August 2024.
  Algorithm~3 (ML-DSA.Verify).

\bibitem{tmldsa2025} A.~Celi, R.~del~Pino, T.~Espitau, G.~Niot,
  T.~Prest, ``Efficient Threshold ML-DSA,'' USENIX Security 2025.

\bibitem{lp020} ``LP-020: Quasar Consensus,'' Lux Industries, 2025.

\bibitem{lp132} ``LP-132: QuasarGPU Adapter,'' Lux Industries, 2025.

\bibitem{hjky97} A.~Herzberg, S.~Jarecki, H.~Krawczyk, M.~Yung,
  ``Proactive Secret Sharing or: How to Cope With Perpetual Leakage,''
  \emph{CRYPTO}, 1995, updated 1997.

\bibitem{desmedtjajodia97} Y.~Desmedt, S.~Jajodia, ``Redistributing
  Secret Shares to New Access Structures and Its Applications,''
  Technical Report ISSE TR-97-01, George Mason University, 1997.

\bibitem{wongwangwing02} T.~M.~Wong, C.~Wang, J.~M.~Wing, ``Verifiable
  Secret Redistribution for Archive Systems,''
  \emph{IEEE SISW}, 2002.

\bibitem{lp073} ``LP-073: Pulsar --- Dynamic Lattice Threshold
  Signatures for Permissionless Validator Sets,'' Lux Industries, 2026.

\end{thebibliography}


%% ==================================================================
\appendix
%% ==================================================================

%% ==================================================================
\section{Threshold ML-DSA Rejection in 3.0}
\label{app:tmldsa}
%% ==================================================================

3.0 final does not adopt threshold ML-DSA.  This appendix records why,
because the trade-off recurs in every audit cycle.

\subsection{The rejection-sampling circularity}

ML-DSA-65 (FIPS~204, Algorithm~2) signs by repeatedly sampling
$\mathbf{y} \leftarrow [-\gamma_1, \gamma_1]^\ell$, computing
$\mathbf{w} = \mathbf{A}\mathbf{y}$, deriving challenge
$c = H(\mu, \mathbf{w}_1)$, then computing $\mathbf{z} =
\mathbf{y} + c \mathbf{s}_1$ and \emph{rejecting} the trial if
$\|\mathbf{z}\|_\infty \geq \gamma_1 - \beta$ or
$\|\text{LowBits}(\mathbf{w} - c\mathbf{s}_2)\|_\infty \geq \gamma_2 -
\beta$.  The rejection probability is non-negligible (roughly $1 - 1/e$
for the $\mathbf{z}$ check), so the signer iterates until accepted.

Threshold ML-DSA must produce a signature whose distribution is
identical to centralised ML-DSA, including the rejection condition
applied to the \emph{aggregate} $\mathbf{z} = \sum_i \mathbf{z}_i$.
Each share $\mathbf{z}_i = \mathbf{y}_i + c \mathbf{s}_{1,i}$ is a
linear function of the per-share randomness $\mathbf{y}_i$.  The
challenge $c$ is a function of the aggregate commitment
$\mathbf{w} = \sum_i \mathbf{w}_i$, which fixes
$\mathbf{y} = \sum_i \mathbf{y}_i$.

The circularity:
\begin{itemize}[nosep]
  \item Each signer must commit $\mathbf{w}_i$ \emph{before} seeing
    $c$ (otherwise it is not a Fiat--Shamir signature).
  \item The acceptance probability of the resulting aggregate
    $\mathbf{z}$ depends on the joint distribution of all
    $\mathbf{y}_i$, which the protocol cannot influence post-commit.
  \item Naive aggregation either (a) leaks information about each
    signer's $\mathbf{s}_{1,i}$ on rejection (because rejection is a
    function of $\mathbf{z}_i$ and a partial $c$), or (b) requires a
    fresh round per rejection, blowing up the round complexity by an
    expected factor of $e$ and breaking the two-round target.
\end{itemize}

The Celi--del~Pino--Espitau--Niot--Prest construction
(USENIX Security 2025~\cite{tmldsa2025}) addresses this with masked
shares and amortised rejection, but the proof is in the algebraic
group model with non-trivial trapdoor assumptions and is not yet
audit-clean.  3.0 freezes on the alternative below.

\subsection{The 3.0 alternative}

Instead of threshold ML-DSA we use:
\begin{itemize}[nosep]
  \item \emph{Per-validator individual} ML-DSA-65 signatures (no
    threshold, no aggregation, no rejection circularity).
  \item A Z-Chain Groth16 rollup proof that all $n$ such signatures
    verify against $\mathsf{zchain\_vk\_root}$ on a common
    $\mathsf{subj}$.
\end{itemize}

This preserves the security target (each validator's signature is
EUF-CMA under MLWE/MSIS, the rollup is knowledge-sound under KEA in the
AGM) without the circularity.  The cost is an $n \times $ blow-up in
witness size, which is precisely what Groth16 amortises across
validators.  See App.~B for the constraint count.

%% ==================================================================
\section{ML-DSA-65 Groth16 Circuit Analysis}
\label{app:circuit}
%% ==================================================================

This appendix derives the R1CS constraint count for verifying a single
ML-DSA-65 signature inside a Groth16 circuit, establishes the
circuit-independent proof size, and corrects the prover-time estimate.

\subsection{ML-DSA-65 Verification Steps}

ML-DSA-65 verification (Algorithm~3 of FIPS~204~\cite{nist-mldsa-spec})
performs the following on input $(pk, \mu, \sigma)$ where
$pk = (\rho, t_1)$ and $\sigma = (\tilde{c}, z, h)$:

\begin{enumerate}[nosep]
  \item Expand $\mathbf{A} \in R_q^{k \times \ell}$ from $\rho$
    via SHAKE-128 ($k=6$, $\ell=5$, $q = 8\,380\,417 = 2^{23} - 2^{13} + 1$).
  \item Compute $c = \mathsf{SampleInBall}(\tilde{c})$: a polynomial in
    $R_q$ with coefficients in $\{-1,0,1\}$, Hamming weight $\tau=49$.
  \item Compute $\mathbf{w}' = \mathbf{A}\mathbf{z} - c \cdot t_1 \cdot 2^d$
    via NTT-domain arithmetic.
  \item Extract high bits: $\mathbf{w}_1' = \mathsf{UseHint}(h, \mathbf{w}')$.
  \item Check $\|z\|_\infty < \gamma_1 - \beta$ where
    $\gamma_1 = 2^{19}$, $\beta = 196$.
  \item Recompute $\tilde{c}' = H(\mu, \mathbf{w}_1')$ and check
    $\tilde{c}' = \tilde{c}$.
\end{enumerate}

\subsection{R1CS Constraint Count}

Each polynomial in $R_q = \mathbb{Z}_q[X]/(X^{256}+1)$ has 256
coefficients modulo $q$; each coefficient fits in 23 bits.

\paragraph{NTT.}
A 256-point NTT over $\mathbb{Z}_q$ requires $256 \cdot \log_2(256) / 2
= 256 \cdot 4 = 1\,024$ butterfly operations. Each butterfly is
one multiplication and two additions in $\mathbb{Z}_q$.  In R1CS,
multiplication modulo $q$ costs 1 constraint (the multiplication gate)
plus $O(1)$ range-check constraints to enforce that the result is reduced
modulo $q$.  Using standard range-check decomposition into 23-bit limbs,
each $\bmod\,q$ reduction costs $\approx 3$ R1CS constraints.  Total
per NTT: $\approx 1\,024 \times 4 = 4\,096$ constraints.

\paragraph{Matrix-vector product $\mathbf{A}\mathbf{z}$.}
This is $k \times \ell = 6 \times 5 = 30$ polynomial multiplications
in $R_q$.  Each polynomial multiplication via NTT requires: forward NTT
of both operands (2 NTTs), pointwise multiply (256 multiplications),
inverse NTT of the result (1 NTT).  However, $\mathbf{A}$ can be
precomputed in NTT domain as a public input, saving $k \times \ell = 30$
forward NTTs.  Remaining: 5 NTTs for $\mathbf{z}$, 30 pointwise
multiplies (each 256 R1CS multiplications + mod-$q$ reductions), 6
inverse NTTs for the $k$ result polynomials, plus 5 additions of
$\ell$-way sums.

\begin{itemize}[nosep]
  \item Forward NTT of $\mathbf{z}$: $5 \times 4\,096 = 20\,480$
  \item Pointwise multiply: $30 \times 256 \times 4 = 30\,720$
  \item Inverse NTT of results: $6 \times 4\,096 = 24\,576$
  \item Additions (free in R1CS): 0
\end{itemize}
Subtotal: $\approx 75\,776$ constraints.

\paragraph{Scalar-polynomial product $c \cdot t_1 \cdot 2^d$.}
The challenge $c$ has weight $\tau = 49$, so $c \cdot t_1$ involves a
sparse convolution.  In the worst case this is done via NTT:
$2 + 1 = 3$ NTTs plus 256 pointwise multiplies.
Cost: $3 \times 4\,096 + 256 \times 4 = 13\,312$ constraints.  Repeated
for each of the $k = 6$ polynomials in $t_1$: $\approx 79\,872$.
(Optimization: exploit sparsity of $c$ to reduce to $\approx 49 \times
256 \times k$ multiply-add operations $\approx 75\,264$.)

\paragraph{Norm check, UseHint, hash.}
Norm check $\|z\|_\infty < \gamma_1 - \beta$: range check on
$5 \times 256 = 1\,280$ coefficients, each $\approx 20$ constraints
for a 19-bit comparison.  Cost: $\approx 25\,600$.
$\mathsf{UseHint}$: bit extraction on $6 \times 256$ coefficients,
$\approx 6 \times 256 \times 10 = 15\,360$ constraints.
$\mathsf{SHAKE}$ / hash re-computation: SHAKE-256 on $\sim$1\,KB
input, $\approx 25\,000$ constraints (Keccak-f[1600] at $\approx
12\,500$ constraints per permutation, 2 calls).

\paragraph{SHAKE-128 for $\mathbf{A}$ expansion.}
Expanding $\mathbf{A}$ from $\rho$ requires $k \times \ell = 30$
rejection-sampling streams of SHAKE-128.  Each stream produces
$\sim 256$ coefficients in $\mathbb{Z}_q$ via rejection sampling.
This is the most expensive part: $\approx 30 \times 3 \times 12\,500
= 1\,125\,000$ constraints (3 Keccak permutations per stream on
average due to rejection).

\paragraph{Total for one ML-DSA-65 verification.}

\begin{center}
\begin{tabular}{lr}
\toprule
\textbf{Sub-circuit} & \textbf{Constraints} \\
\midrule
$\mathbf{A}$ expansion (SHAKE-128)    & $\sim 1\,125\,000$ \\
$\mathbf{A}\mathbf{z}$ product (NTT)  & $\sim 75\,776$ \\
$c \cdot t_1$ product                  & $\sim 79\,872$ \\
Norm check                             & $\sim 25\,600$ \\
UseHint                                & $\sim 15\,360$ \\
Hash recomputation (SHAKE-256)         & $\sim 25\,000$ \\
\midrule
\textbf{Total (1 verification)}        & $\sim 1.35 \times 10^6 \approx 2^{22.5}$ \\
\bottomrule
\end{tabular}
\end{center}

For $n$ validators, the circuit verifies $n$ signatures:
$C \approx n \times 1.35 \times 10^6$.  For $n = 21$ (Quasar mainnet):
$C \approx 2.84 \times 10^7 \approx 2^{24.8}$.

\paragraph{Optimisation -- amortised case.}
If $\mathbf{A}$ is shared across all $n$ validators (same $\rho$),
the expansion is done once: $C \approx 1.125 \times 10^6 + n \times
2.21 \times 10^5 \approx 5.77 \times 10^6 \approx 2^{22.5}$ for
$n = 21$.  This is the figure used in the main text and in LP-020
\S6.4.  Equivalently, the per-validator amortised constraint count is
$\sim 2^{20}$.

\subsection{Proof Size}

Groth16 proof size is circuit-independent~\cite{groth2016}: two
elements of $\mathbb{G}_1$ and one element of $\mathbb{G}_2$ over
BLS12-381:
\[
  |\pi| = 2 \times 48 + 1 \times 96 = 192 \text{ bytes}
\]
This holds regardless of the constraint count $C$.  The 192-byte claim
in the main text (Def.~6.2 of LP-020) is correct.

\subsection{Prover Time}

Groth16 proving requires two multi-scalar exponentiations of size $C$
and an FFT of size $C$.  The dominant cost is the MSM:
$O(C / \log C)$ group operations in $\mathbb{G}_1$.  On BLS12-381,
one $\mathbb{G}_1$ scalar multiplication takes $\sim 0.5$\,ms on a
single CPU core (Apple M1 class).

For $C \approx 2^{22.5}$ (optimised, $n = 21$):
\begin{itemize}[nosep]
  \item MSM via Pippenger: $C / \log_2 C \approx 2^{22.5} / 22.5
    \approx 2.6 \times 10^5$ group additions.
  \item At $\sim 1\,\mu$s per $\mathbb{G}_1$ addition: $\approx 0.26$\,s.
  \item FFT: $O(C \log C)$ field operations, $\approx 0.1$\,s.
  \item \textbf{Total prover time: $\approx 0.4$\,s on a single CPU core.}
  \item With GPU acceleration (e.g., \texttt{bellman} MSM on RTX 4090):
    $\approx 5\text{--}15$\,ms.
\end{itemize}

\paragraph{Correction to earlier drafts.}
The $357\,\mu$s figure cited in earlier drafts was unsupported by
measurement.  The closest real numbers are BLS single signing
($350\,\mu$s) and ML-DSA-65 signing ($495\,\mu$s).  For the actual
ML-DSA-65 verification circuit at $C \approx 2^{22.5}$, CPU prover time
is $\sim 400$\,ms and GPU prover time is $\sim 5\text{--}15$\,ms.  Full
QuasarCert CPU verification (BLS aggregated verify + Groth16 verify) is
$\sim 2\text{--}5$\,ms.  Z-Chain proof generation runs asynchronously
(epoch-based, not per-block), so this does not impact finality
latency.  See \texttt{consensus/BENCHMARKS.md} for all measured values.


%% ==================================================================
\section{Static vs Adaptive Corruption}
\label{app:adaptive}
%% ==================================================================

Theorem~\ref{thm:main} and its counterpart Theorem~7.5 in LP-020
assume \emph{static} corruption: the adversary commits to the set of
corrupted validators before protocol execution begins.  This appendix
addresses extension to the adaptive corruption model.

\subsection{Corruption Models}

\begin{definition}[Static corruption~\cite{canetti2001}]
The adversary $\mathcal{A}$ selects the set $F \subset \mathcal{V}$
of corrupted parties before the protocol begins.  $\mathcal{A}$ learns
all internal state of parties in $F$ but cannot corrupt additional
parties during execution.
\end{definition}

\begin{definition}[Adaptive corruption~\cite{canetti2001}]
The adversary $\mathcal{A}$ may corrupt parties at any point during
protocol execution.  Upon corrupting $v_i$, $\mathcal{A}$ learns all
of $v_i$'s current state (keys, randomness, signing transcripts).
The total stake of corrupted parties remains bounded by $\tau$.
\end{definition}

In the adaptive model, the simulator in a security proof must produce
a consistent view for the adversary \emph{after} the adversary chooses
which party to corrupt.  This is strictly harder than the static model:
the simulator must either (a)~program the random oracle to maintain
consistency (programmable RO), or (b)~ensure honest parties have
erased sensitive state before corruption (secure erasures).

\subsection{Static-Model Limitation}

\begin{remark}
Theorem~\ref{thm:main} holds only under static corruption.  The proofs
of Ringtail SUF-CMA~\cite{ringtail} and BLS
EUF-CMA~\cite{boneh2001short} both use rewinding-based extractors
that require the corrupted set to be fixed before the simulation
begins.  In the adaptive model, the adversary can corrupt a validator
\emph{after} seeing its Ringtail Round~1 commitment, learning the
signing randomness $\mathbf{y}_i$ and secret share $\mathbf{s}_i$.
\end{remark}

\subsection{Extension via Programmable Random Oracle}

\begin{corollary}[Adaptive soundness, programmable RO]
\label{cor:adaptive-pro}
In the programmable random oracle model, Theorem~\ref{thm:main} extends
to adaptive corruption with a security loss of $O(n \cdot Q_H)$, where
$Q_H$ is the number of hash queries.
\end{corollary}

\begin{proof}[Proof sketch]
We apply the Fischlin transform~\cite{fischlin2005}: the simulator
programs the random oracle $H$ so that upon adaptive corruption of
$v_i$, the simulated signing transcripts remain consistent with the
adversary's view.  Per-component: BLS uses RO programming to embed a
co-CDH challenge at the corruption point (loss $O(n)$); Ringtail
programs $H$ on the round-1 commitment so the adversary's view is
extractable, then solves a linear system over $R_q$ on corruption
(loss $O(n)$); Groth16 in the AGM is rewinding-free and adaptively
sound directly.

Combining with a union bound over $n$ targets and $Q_H$ queries:
\[
  \mathsf{Adv}^{\text{forge}}_{\text{adaptive}} \leq
  n \cdot Q_H \cdot \max(\epsilon_{\text{co-CDH}},
  \epsilon_{\text{MLWE}}, \epsilon_{\text{MLWE/MSIS}})
\]
which is negligible for $\lambda = 128$ when $n \cdot Q_H < 2^{60}$
(a generous bound for any practical deployment).
\end{proof}

\subsection{Extension via Secure Erasures}

\begin{corollary}[Adaptive soundness, erasure model]
\label{cor:adaptive-erasure}
If honest validators erase signing randomness
($\mathbf{y}_i$, PRF state) after each Ringtail round completes,
then Theorem~\ref{thm:main} extends to adaptive corruption without
the programmable RO assumption, with no additional security loss.
\end{corollary}

\begin{proof}[Proof sketch]
If $v_i$ erases $\mathbf{y}_i$ after broadcasting $\mathbf{z}_i$,
then adaptive corruption of $v_i$ reveals only $\mathbf{s}_i$ (the
long-term share) and the completed signature transcript.  The
adversary gains no advantage beyond what a static adversary controlling
$v_i$ from the start would have.  The static-model proof applies
directly.
\end{proof}

\paragraph{Practical note.}
The Quasar implementation (Section~9 of LP-020) zeroes signing
randomness after each round via Go's explicit memory clearing.
This provides the erasure property in software.  HSM-backed signing
(Section~5 of LP-020, T-Chain) provides hardware-enforced erasure.

\paragraph{Blockchain corruption model.}
In blockchain settings, the standard corruption model is that of
Garay--Kiayias--Leonardos~\cite{garay2015}: the adversary is
\emph{mildly adaptive}, corrupting parties between rounds but not
within a round.  Both Cor.~\ref{cor:adaptive-pro} and
Cor.~\ref{cor:adaptive-erasure} cover this model as a special
case.


%% ==================================================================
\section{Trusted-Setup Ceremony}
\label{app:setup}
%% ==================================================================

\subsection{The SRS Assumption}

Groth16~\cite{groth2016} requires a structured reference string (SRS)
of the form:
\[
  \mathsf{SRS} = \bigl(
    \{g^{\tau^i}\}_{i=0}^{C},\;
    \{h^{\tau^i}\}_{i=0}^{C},\;
    \{g^{\alpha \tau^i}\},\;
    \{g^{\beta \tau^i}\},\;
    h^\beta,\; h^\gamma,\; h^\delta
  \bigr)
\]
where $\tau, \alpha, \beta, \gamma, \delta$ are toxic waste that must
be destroyed after generation.  If any party learns $\tau$, it can
forge proofs for arbitrary false statements.

\paragraph{Circuit specificity.}
The Groth16 SRS is circuit-specific: a new SRS must be generated for
each distinct R1CS circuit.  For QuasarCert, the circuit is fixed
(ML-DSA-65 verification for a given $n$), so a single ceremony per
$(n, \text{ML-DSA parameters})$ tuple suffices.

\subsection{Multi-Party Ceremony}

The standard mitigation is an $N$-party ceremony where each
participant $P_j$ contributes randomness $\tau_j$ to the SRS.  The
final SRS uses $\tau = \sum_j \tau_j$ (or $\prod_j \tau_j$ in the
multiplicative formulation).  Security holds if \emph{at least one}
participant is honest and destroys their $\tau_j$.

\begin{theorem}[Ceremony security~\cite{bowe2017}]
An $N$-party Groth16 SRS ceremony is secure (no party learns $\tau$)
provided at least one participant honestly destroys their contribution.
The ceremony can be run sequentially (each participant extends the
previous SRS) with public verification of each step.
\end{theorem}

For QuasarCert with $n = 21$ validators, the ceremony participants can
be the validators themselves.  With 21 participants, the probability
that all are dishonest (and collude to reconstruct $\tau$) is
negligible under the $f < 1/3$ Byzantine assumption: at most 7
validators are corrupt, so at least 14 honestly destroy their
contributions.  The output is committed in
$\mathsf{zchain\_vk\_root}$ on the Z-Chain genesis block.

\subsection{Universal-Setup Alternatives}

Universal-setup proof systems avoid circuit-specific ceremonies at the
cost of larger proofs.

\begin{center}
\begin{tabular}{lrrll}
\toprule
\textbf{System} & \textbf{Proof size} & \textbf{Verify} &
  \textbf{Setup} & \textbf{Ref.} \\
\midrule
Groth16 & 192\,B & 3 pairings & Circuit-specific & \cite{groth2016} \\
PLONK   & $\sim$500\,B & $O(1)$ pairings & Universal, updatable
  & \cite{plonk2019} \\
Marlin  & $\sim$800\,B & $O(\log C)$ & Universal, updatable
  & \cite{marlin2020} \\
\bottomrule
\end{tabular}
\end{center}

\paragraph{Recommendation (3.0 final).}
Groth16 with a multi-party ceremony among the validator set, output
committed in $\mathsf{zchain\_vk\_root}$, is the production
configuration for the 2025-12-25 activation.  PLONK migration is
tracked as a future-work item under LP-020 \S11.


%% ==================================================================
\section{Quantum-Safety Analysis}
\label{app:quantum}
%% ==================================================================

This appendix consolidates the concrete classical and quantum security
estimates for each lane, then states the combined three-lane survival
condition.

\subsection{Lane 1: BLS12-381 (classical only)}

Classical security is approximately $2^{128}$ (best classical attack
on the discrete log in $\mathbb{F}_{q^{12}}$ via the special tower
NFS).  Quantum security is $0$ in polynomial time once a
cryptographically relevant quantum computer exists, by Shor on the
underlying $\mathbb{G}_1, \mathbb{G}_2$.

\subsection{Lane 2: Ringtail (Module-LWE / Module-SIS)}

Parameter set from~\cite{ringtail}:

\begin{center}
\begin{tabular}{lll}
\toprule
\textbf{Parameter} & \textbf{Symbol} & \textbf{Value} \\
\midrule
Ring dimension   & $N$         & 256 \\
Modulus          & $q$         & $2^{48} + 2561$ ($\approx 2^{48}$) \\
Module rows      & $M$         & 8 \\
Module columns   & $N'$        & 7 \\
Gaussian width (keygen) & $\sigma_e$ & 6.108 \\
Gaussian width (signing)& $\sigma^*$ & $2^{37.33}$ \\
Rejection bound  & $B$         & $2^{48.61}$ \\
Challenge weight & $\kappa$    & 23 \\
\bottomrule
\end{tabular}
\end{center}

The Module-LWE instance underlying Ringtail has effective lattice
dimension $d = M \times N = 8 \times 256 = 2\,048$ over
$\mathbb{Z}_q$.

\paragraph{Classical security via primal attack.}
Best classical attack: the primal lattice attack with BKZ-$\beta$
sieving (BDGL~\cite{bdgl2016}).  Solving requires
$\delta^{2\beta - d - 1} q^{(d+1)/\beta} = \sigma_e \sqrt{2\pi e /
(2\beta)}$ with root Hermite factor
$\delta = \bigl((\pi\beta)^{1/\beta} \beta / (2\pi e)\bigr)^{1/(2(\beta
-1))}$.  For $d = 2048, q \approx 2^{48}, \sigma_e = 6.108$, numerics
give $\beta \approx 430$.  Classical cost
\[
  T_{\mathrm{classical}} = 2^{0.292\beta + c_0}
  \approx 2^{0.292 \cdot 430 + 16.4}
  = 2^{142}.
\]

\paragraph{Quantum security via Grover-sieving.}
The Laarhoven--BDGL exponent under Grover speedup is $0.265$:
\[
  T_{\mathrm{quantum}} = 2^{0.265 \cdot 430 + 16.4} = 2^{130.35}.
\]

\subsection{Lane 3a: ML-DSA-65 (FIPS~204)}

Parameters: $k = 6$, $\ell = 5$, $q = 2^{23} - 2^{13} + 1$, $\eta = 4$.
Effective MLWE dimension $5 \cdot 256 = 1280$, $\log_2 q = 23$.
NIST Cat.~3 target: $\geq 2^{128}$ classical / $\geq 2^{128}$ quantum
under the same NIST cost model~\cite{nist-pqc-cost}.  Best known attack
gives $\beta_{\mathrm{BKZ}} \approx 360$, $T_{\mathrm{quantum}} \geq
2^{128}$.

\subsection{Lane 3b: Z-Chain Groth16}

Classical: $2^{128}$ (BLS12-381 pairing-based, AGM/KEA).  Quantum:
Groth16 itself loses classical knowledge-soundness once Shor breaks
discrete log on $\mathbb{G}_1$.  However, the \emph{statement proven}
remains a conjunction of $n$ ML-DSA-65 signature verifications, so a
quantum adversary's effective task degrades to forging ML-DSA-65
signatures committed by $\mathsf{zchain\_vk\_root}$.  Standalone PQ
security of Lane~3 thus tracks Lane~3a above.

A future-work direction is to swap Groth16 for a STARK-of-pairing or
hash-based SNARK that is itself PQ-knowledge-sound.  Tracked in
LP-020 \S11; not in 3.0 final scope.

\subsection{Comparison table}

\begin{center}
\begin{tabular}{lrrrr}
\toprule
\textbf{Scheme} & $d$ & $\log_2 q$ & $\beta_{\text{BKZ}}$ &
  $T_{\text{quantum}}$ \\
\midrule
ML-DSA-65   & 1280 & 23 & $\approx 360$
  & $\geq 2^{128}$ \\
ML-KEM-768  & 768  & 12 & $\approx 340$
  & $\geq 2^{128}$ \\
Ringtail    & 2048 & 48 & $\approx 430$
  & $\geq 2^{130}$ \\
BLS12-381   & --   & -- & --
  & $0$ (Shor) \\
\bottomrule
\end{tabular}
\end{center}

\subsection{Combined survival condition}

The full QuasarCert quantum-safe predicate is
$\mathsf{VerifyPQ} = \mathsf{RT.Verify} \land \mathsf{ZK.Verify}$
(Theorem~\ref{thm:pq}).  Equivalently, ``1 of 2'' PQ pillars must hold:
either Ringtail's MLWE/MSIS, or the ML-DSA-65 path inside the Groth16
proof.  The classical BLS lane provides classical-only defense in
depth and contributes nothing to the PQ guarantee, in line with the
analysis of pq-finality-no-bls.tex.

\subsection{Caveats}

\begin{enumerate}[nosep]
  \item BDGL sieving exponents $0.292$ (classical) and $0.265$
    (quantum) are heuristic.
  \item The Grover model assumes ideal oracle access; physical
    overhead increases effective security.
  \item Ring structure of $\mathbb{Z}_q[X]/(X^{256}+1)$ has resisted
    known algebraic attacks at the relevant dimensions; an algorithmic
    breakthrough would require parameter revision.
  \item Lane~3b PQ guarantee is conditional on the future-work SNARK
    swap.  In 3.0 final, Lane~3 is treated as classical with PQ
    fallback to its statement.
\end{enumerate}


\end{document}
